%!TEX program = lualatex
% Relatório: Experimento 1 — Medidas Elétricas

\documentclass[a4paper,12pt]{article}

% =========================================================
% Idioma e tipografia
% =========================================================
\usepackage[provide=*,brazilian]{babel}
\usepackage{iftex}
\ifLuaTeX
  \usepackage{fontspec}
  \setmainfont{Latin Modern Roman}
\else
  \usepackage[T1]{fontenc}
  \usepackage[utf8]{inputenc}
\fi

% =========================================================
% Layout
% =========================================================
\usepackage[top=3cm,bottom=2.5cm,left=3cm,right=2.5cm]{geometry}
\usepackage{setspace}
\onehalfspacing
\usepackage{indentfirst}
\setlength{\parindent}{1.25cm}

% =========================================================
% Matemática, tabelas, listas e figuras
% =========================================================
\usepackage{amsmath,amssymb}
\usepackage{float}
\usepackage{booktabs}
\usepackage{multirow}
\usepackage{array}
\usepackage{enumitem}
\usepackage{graphicx}

\setlist[itemize]{noitemsep,topsep=3pt}
\setlist[enumerate]{noitemsep,topsep=3pt}

\graphicspath{{./}{./Figs/}{./imagens/}{./figuras/}}

% Mantém o documento compilável mesmo quando alguma imagem ainda não está na pasta.
\newcommand{\imagemrelatorio}[2][width=.5\linewidth]{%
  \IfFileExists{#2}{%
    \includegraphics[#1]{#2}%
  }{%
    \fbox{\parbox[c][3.5cm][c]{.78\linewidth}{%
      \centering Imagem não encontrada: \texttt{\detokenize{#2}}%
    }}%
  }%
}

% =========================================================
% Links
% =========================================================
\usepackage{xurl}
\usepackage{hyperref}
\hypersetup{
  colorlinks=true,
  linkcolor=black,
  urlcolor=blue,
  citecolor=black,
  pdfauthor={José Mauro Elsas},
  pdftitle={Experimento 1 — Medidas Elétricas}
}

% =========================================================
% Dados da capa
% =========================================================
\newcommand{\Instituicao}{Universidade do Estado do Rio de Janeiro}
\newcommand{\Curso}{Física VII -- Oceanografia}
\newcommand{\Professor}{Professor: Alexandre Magnus Gomes Carvalho}
\newcommand{\Aluno}{Aluno: José Mauro Elsas}
\newcommand{\Cidade}{Rio de Janeiro}
\newcommand{\Ano}{2026}

\begin{document}
\pagenumbering{gobble}

% =========================================================
% Capa
% =========================================================
\begin{titlepage}
  \centering
  {\large \Instituicao\par}
  \vspace{0.5cm}
  {\large \Curso\par}
  \vspace{1cm}

  {\LARGE\textbf{Experimento 1}\par}
  \vspace{0.6cm}
  {\Large Medidas Elétricas\par}

  \vspace{1cm}
  \imagemrelatorio[width=.3\linewidth]{uerj.jpg}

  \vfill

  \begin{flushright}
    \Aluno\\
    \Professor
  \end{flushright}

  \vspace{1cm}
  {\Cidade{} --- \Ano\par}
\end{titlepage}
\clearpage
\pagenumbering{arabic}

% =========================================================
% Objetivos
% =========================================================
\section{Objetivos}

São os seguintes os objetivos deste relatório:

\begin{enumerate}[label=\alph*.]
  \item Determinar os desvios-padrão das medições da fonte elétrica, tanto em
  tensão contínua quanto alternada;

  \item Determinar os desvios-padrão dos valores medidos das resistências e
  verificar se o valor de referência de cada resistência está no intervalo de
  medição.
\end{enumerate}

% =========================================================
% Introdução
% =========================================================
\section{Introdução}

O estudo da eletricidade desenvolveu-se a partir da observação de fenômenos
naturais e de experimentos realizados ao longo de vários séculos. Embora
manifestações elétricas fossem conhecidas desde a Antiguidade, foi principalmente
entre os séculos XVIII e XIX que a eletricidade passou a ser investigada de
maneira sistemática. Posteriormente, a necessidade de comparar resultados levou
ao desenvolvimento e à padronização internacional de unidades de tensão,
corrente e resistência \cite{nist_ampere}.

Com o avanço da instrumentação elétrica, tornou-se necessário avaliar não apenas
o valor indicado por um aparelho, mas também a qualidade e a confiabilidade da
medição. Nenhuma medição experimental é absolutamente exata, pois seus resultados
podem ser influenciados pela resolução e pela calibração dos instrumentos, pelas
condições ambientais, pelas características do circuito e pelo procedimento
adotado. Assim, medições repetidas de uma mesma grandeza podem apresentar pequenas
diferenças, formando um conjunto de valores distribuídos em torno de um valor
médio.

A análise estatística permite descrever essa dispersão e avaliar a repetibilidade
dos resultados. A média aritmética representa uma estimativa do valor central das
medições, enquanto o desvio-padrão caracteriza sua dispersão. Uma dispersão menor
indica maior precisão das medições, mas não garante, isoladamente, sua exatidão,
pois os resultados podem estar concentrados em torno de um valor diferente do
valor de referência \cite{jcgm_vim,jcgm_gum}.

Neste experimento, as medições de tensão contínua e de tensão alternada são
analisadas separadamente, pois correspondem a diferentes modos de funcionamento e
de leitura do instrumento. Para as resistências elétricas, além da determinação
da média e do desvio-padrão, comparam-se os resultados experimentais com os
respectivos valores de referência. Dessa forma, os conceitos de média,
desvio-padrão, precisão, exatidão e intervalo de medição fornecem a base teórica
para avaliar a dispersão das medições da fonte elétrica e das resistências.

% =========================================================
% Materiais e métodos
% =========================================================
\section{Materiais e Métodos}

Neste experimento foram utilizados os seguintes materiais:

\begin{enumerate}[label=\alph*.]
  \item Multímetro digital com as respectivas pontas de prova.

  \begin{figure}[H]
    \centering
    \imagemrelatorio{Figs/digital.jpg}
    \caption{Multímetro digital.}
    \label{fig:multimetro_digital}
  \end{figure}

  \item Fonte de alimentação CA/CC estabilizada.

  \begin{figure}[H]
    \centering
    \imagemrelatorio{Figs/fonte.jpeg}
    \caption{Fonte de alimentação CA/CC.}
    \label{fig:fonte_alimentacao}
  \end{figure}

  \item \textit{Protoboard}, placa de ensaio constituída por uma base plástica
  com orifícios e conexões internas condutoras, utilizada para montar e testar
  circuitos eletrônicos sem a necessidade de solda \cite{ufg_protoboard}.

  \begin{figure}[H]
    \centering
    \imagemrelatorio{Figs/protoboard.jpeg}
    \caption{\textit{Protoboard}.}
    \label{fig:protoboard}
  \end{figure}

  \item Resistores de diferentes valores nominais.

  \begin{figure}[H]
    \centering
    \imagemrelatorio{Figs/resistores.jpeg}
    \caption{Resistores utilizados no experimento.}
    \label{fig:resistores}
  \end{figure}
\end{enumerate}

\subsection{Procedimento experimental}

Para medir as tensões contínua ($V_{\mathrm{CC}}$) e alternada
($V_{\mathrm{CA}}$), a fonte de alimentação foi conectada à rede elétrica e
ligada, mantendo-se seus controles de tensão e corrente em condições adequadas
para as medições. Em seguida, o seletor de ajuste da fonte foi posicionado
sucessivamente nas marcações 1, 4, 7 e 10.

Em cada posição foram efetuadas cinco medições de tensão contínua e cinco de
tensão alternada. As tensões foram medidas com o multímetro conectado em paralelo
aos terminais de saída da fonte. Posteriormente, os resistores foram montados na
\textit{protoboard} e suas resistências elétricas foram medidas com o multímetro
digital na função de ohmímetro.

Para cada conjunto de cinco leituras, calculou-se a média aritmética:

\[
  \bar{x}=\frac{1}{n}\sum_{i=1}^{n}x_i,
\]

em que $x_i$ representa cada valor medido e $n=5$ é o número de medições em
cada posição. O desvio-padrão amostral foi calculado por:

\[
  s=\sqrt{\frac{\sum_{i=1}^{n}(x_i-\bar{x})^2}{n-1}}.
\]

Os resultados das resistências podem ser apresentados na forma

\[
  R=(\bar{R}\pm s)\,\Omega,
\]

em que $\bar{R}$ é a resistência média e $s$ é o desvio-padrão amostral.

% =========================================================
% Resultados
% =========================================================
\section{Resultados}

\subsection{Medições de tensão}

\begin{table}[H]
  \centering
  \caption{Valores de tensão observados nas diferentes posições da fonte.}
  \label{tab:medicoes_tensao}
  \medskip

  \begin{tabular}{ccc}
    \toprule
    {\textbf{Posição}}
    & {\textbf{Tensão contínua} ($\mathrm{V}$)}
    & {\textbf{Tensão alternada} ($\mathrm{V}$)} \\
    \midrule
    \multirow{5}{*}{1}
      & 2,7  & 2,0 \\
      & 2,6  & 2,0 \\
      & 2,6  & 2,0 \\
      & 2,6  & 2,0 \\
      & 2,6  & 2,0 \\
    \midrule
    \multirow{5}{*}{4}
      & 11,5 & 8,5 \\
      & 11,5 & 8,6 \\
      & 11,5 & 8,6 \\
      & 11,5 & 8,6 \\
      & 11,5 & 8,6 \\
    \midrule
    \multirow{5}{*}{7}
      & 20,5 & 14,8 \\
      & 20,5 & 14,8 \\
      & 20,5 & 14,8 \\
      & 20,5 & 14,8 \\
      & 20,5 & 14,8 \\
    \midrule
    \multirow{5}{*}{10}
      & 29,5 & 21,4 \\
      & 29,5 & 21,4 \\
      & 29,5 & 21,4 \\
      & 29,5 & 21,4 \\
      & 29,5 & 21,3 \\
    \bottomrule
  \end{tabular}
\end{table}

\subsection{Médias das leituras de tensão}

\begin{table}[H]
    \centering
    \caption{Valores médios e desvios-padrão amostrais das tensões
    contínua e alternada.}
    \label{tab:medias_desvios_vcc_vca}
    \medskip

    \begin{tabular}{ccccc}
        \toprule
        \textbf{Posição}
        & \shortstack{\textbf{Média de}\\
          \boldmath$V_{\mathrm{CC}}$\unboldmath\\($\mathrm{V}$)}
        & \shortstack{\textbf{Desvio-padrão}\\
          \boldmath$s_{\mathrm{CC}}$\unboldmath\\($\mathrm{V}$)}
        & \shortstack{\textbf{Média de}\\
          \boldmath$V_{\mathrm{CA}}$\unboldmath\\($\mathrm{V}$)}
        & \shortstack{\textbf{Desvio-padrão}\\
          \boldmath$s_{\mathrm{CA}}$\unboldmath\\($\mathrm{V}$)}
        \\
        \midrule
        1  & 2,62 & 0,045 & 2,00  & 0     \\
        4  & 11,5 & 0     & 8,58  & 0,045 \\
        7  & 20,5 & 0     & 14,8  & 0     \\
        10 & 29,5 & 0     & 21,38 & 0,045 \\
        \bottomrule
    \end{tabular}
\end{table}

\noindent
\textit{Nota:} quando todas as leituras de uma determinada posição são
idênticas, não há dispersão entre os valores medidos e, consequentemente,
o desvio-padrão amostral é igual a zero (\(s=0\)).

\begin{table}[H]
    \centering
    \caption{Resultados experimentais das medições de tensão contínua
    e alternada.}
    \label{tab:resultados_tensoes}
    \medskip

    \begin{tabular}{ccc}
        \toprule
        \textbf{Posição}
        & \shortstack{\textbf{Tensão contínua}\\
          \(V_{\mathrm{CC}}=(\bar V_{\mathrm{CC}}\pm u)\,\mathrm{V}\)}
        & \shortstack{\textbf{Tensão alternada}\\
          \(V_{\mathrm{CA}}=(\bar V_{\mathrm{CA}}\pm u)\,\mathrm{V}\)}
        \\
        \midrule
        1  & \(2{,}62\pm0{,}04\)  & \(2{,}00\pm0{,}01\)  \\
        4  & \(11{,}50\pm0{,}01\) & \(8{,}58\pm0{,}04\)  \\
        7  & \(20{,}50\pm0{,}01\) & \(14{,}80\pm0{,}01\) \\
        10 & \(29{,}50\pm0{,}01\) & \(21{,}38\pm0{,}04\) \\
        \bottomrule
    \end{tabular}
\end{table}

Embora algumas séries tenham apresentado desvio-padrão igual a zero,
a incerteza das medições não foi considerada nula. Nesses casos,
adotou-se como incerteza instrumental a resolução do multímetro digital,
igual a \(0{,}01\,\mathrm{V}\).

\subsection{Medições das resistências}

\begin{table}[H]
    \centering
    \caption{Valores nominais e leituras das resistências elétricas.}
    \label{tab:leituras_resistores}
    \medskip

    \begin{tabular}{cccc}
        \toprule
        \textbf{Resistor}
        & \shortstack{\textbf{Valor nominal}\\($\Omega$)}
        & \textbf{Leitura}
        & \shortstack{\textbf{Valor medido}\\($\Omega$)}
        \\
        \midrule

        \multirow{3}{*}{R1}
        & \multirow{3}{*}{47}
        & 1   & 48,6 \\
        & & 2 & 48,2 \\
        & & 3 & 48,1 \\
        \midrule

        \multirow{3}{*}{R2}
        & \multirow{3}{*}{100}
        & 1 & 98,9 \\
        & & 2 & 98,8 \\
        & & 3 & 99,0 \\
        \midrule

        \multirow{3}{*}{R3}
        & \multirow{3}{*}{150}
        & 1 & 151,0 \\
        & & 2 & 151,2 \\
        & & 3 & 151,7 \\
        \midrule

        \multirow{3}{*}{R4}
        & \multirow{3}{*}{330}
        & 1 & 323 \\
        & & 2 & 327 \\
        & & 3 & 318 \\
        \bottomrule
    \end{tabular}
\end{table}


\subsection{Médias e desvios-padrão das resistências}

\begin{table}[H]
    \centering
    \caption{Valores nominais, médias e desvios-padrão amostrais
    das resistências elétricas.}
    \label{tab:medias_desvios_resistores}
    \medskip

    \begin{tabular}{cccc}
        \toprule
        \textbf{Resistor}
        & \shortstack{\textbf{Valor nominal}\\($\Omega$)}
        & \shortstack{\textbf{Média medida}\\($\Omega$)}
        & \shortstack{\textbf{Desvio-padrão}\\amostral ($\Omega$)}
        \\
        \midrule
        R1 & 47  & 48,30  & 0,26 \\
        R2 & 100 & 98,90  & 0,10 \\
        R3 & 150 & 151,30 & 0,36 \\
        R4 & 330 & 322,67 & 4,51 \\
        \bottomrule
    \end{tabular}
\end{table}
Os resultados da Tabela~\ref{tab:medias_desvios_resistores}
mostram que o resistor R2 apresentou a menor dispersão entre as três
leituras, com desvio-padrão amostral de \(0,10\,\Omega\). O resistor R4
apresentou a maior dispersão, com desvio-padrão de \(4,51\,\Omega\).
Embora esse valor seja numericamente maior, sua interpretação deve
considerar também a maior resistência nominal desse componente.


A diferença percentual entre o valor medido e o valor nominal foi calculada por

\[
  \varepsilon_{\%}
  =\frac{\lvert R_{\mathrm{medido}}-R_{\mathrm{nominal}}\rvert}
  {R_{\mathrm{nominal}}}\times 100.
\]

Considerando a tolerância nominal de até $5\%$, os resultados podem
ser comparados com o limite especificado pelo fabricante.

\subsection{Resultados das medições das resistências}

Os resultados experimentais das resistências foram expressos na forma

\[
    R=(\bar{R}\pm s)\,\Omega,
\]

em que \(\bar{R}\) representa a média das três leituras e \(s\)
representa o respectivo desvio-padrão amostral.

A diferença percentual entre a resistência média medida e o valor
nominal foi calculada por

\[
    \varepsilon_{\%}
    =
    \frac{
        \left|\bar{R}-R_{\mathrm{nominal}}\right|
    }{
        R_{\mathrm{nominal}}
    }
    \times 100.
\]

\begin{table}[H]
    \centering
    \caption{Resultados experimentais das resistências e comparação
    com os respectivos valores nominais.}
    \label{tab:resultados_resistencias}
    \medskip

    \begin{tabular}{ccccc}
        \toprule
        \textbf{Resistor}
        & \shortstack{\textbf{Valor nominal}\\($\Omega$)}
        & \shortstack{\textbf{Resultado experimental}\\
          \(R=(\bar{R}\pm s)\,\Omega\)}
        & \shortstack{\textbf{Diferença}\\percentual}
        & \shortstack{\textbf{Dentro da}\\tolerância de 5\%?}
        \\
        \midrule
        R1 & 47  & \(48{,}30\pm0{,}26\)  & 2,77\% & Sim \\
        R2 & 100 & \(98{,}90\pm0{,}10\)  & 1,10\% & Sim \\
        R3 & 150 & \(151{,}30\pm0{,}36\) & 0,87\% & Sim \\
        R4 & 330 & \(322{,}67\pm4{,}51\) & 2,22\% & Sim \\
        \bottomrule
    \end{tabular}
\end{table}

\newpage
% =========================================================
% Referências
% =========================================================
\begin{thebibliography}{99}

\bibitem{jcgm_gum}
JOINT COMMITTEE FOR GUIDES IN METROLOGY.
\textit{Evaluation of measurement data: Guide to the expression of uncertainty
in measurement}. JCGM 100:2008. Sèvres: BIPM, 2008.
Disponível em: \url{https://www.bipm.org/en/committees/jc/jcgm/publications}.
Acesso em: 14 ago. 2026.

\bibitem{jcgm_vim}
JOINT COMMITTEE FOR GUIDES IN METROLOGY.
\textit{International vocabulary of metrology: Basic and general concepts and
associated terms}. 3. ed. JCGM 200:2012. Sèvres: BIPM, 2012.
Disponível em: \url{https://www.bipm.org/en/doi/10.59161/jcgm200-2012}.
Acesso em: 14 ago. 2026.

\bibitem{nist_ampere}
NATIONAL INSTITUTE OF STANDARDS AND TECHNOLOGY.
\textit{Ampere: History}. Gaithersburg: NIST, 2018.
Disponível em: \url{https://www.nist.gov/si-redefinition/ampere/ampere-history}.
Acesso em: 14 ago. 2026.

\bibitem{ufg_protoboard}
UNIVERSIDADE FEDERAL DE GOIÁS.
\textit{Protoboard: o que é e como usar?}
Disponível em: \url{https://ipelab.ufg.br/n/156373-protoboard-o-que-e-e-como-usar}.
Acesso em: 14 ago. 2026.

\end{thebibliography}

\end{document}
